\documentclass[10pt,a4paper]{amsart}
\usepackage{amsmath,amssymb,amsthm,graphicx}
\usepackage{verbatim}
\newenvironment{reference}{\comment}{\endcomment}
%\newenvironment{reference}{}{}
\newenvironment{slogan}{\comment}{\endcomment}
\newenvironment{history}{\comment}{\endcomment}

% For commutative diagrams we use Xy-pic
\usepackage[all]{xy}

% We use 2cell for 2-commutative diagrams.
\xyoption{2cell}
% Theorem environments.
%
\theoremstyle{plain}
\newtheorem{theorem}[subsection]{Theorem}
\newtheorem{proposition}[subsection]{Proposition}
\newtheorem{lemma}[subsection]{Lemma}

\theoremstyle{definition}
\newtheorem{definition}[subsection]{Definition}
\newtheorem{example}[subsection]{Example}
\newtheorem{exercise}[subsection]{Exercise}
\newtheorem{situation}[subsection]{Situation}

\theoremstyle{remark}
\newtheorem{remark}[subsection]{Remark}
\newtheorem{remarks}[subsection]{Remarks}

\numberwithin{equation}{subsection}

% Macros
%
\def\lim{\mathop{\mathrm{lim}}\nolimits}
\def\colim{\mathop{\mathrm{colim}}\nolimits}
\def\Spec{\mathop{\mathrm{Spec}}}
\def\Hom{\mathop{\mathrm{Hom}}\nolimits}
\def\Ext{\mathop{\mathrm{Ext}}\nolimits}
\def\SheafHom{\mathop{\mathcal{H}\!\mathit{om}}\nolimits}
\def\SheafExt{\mathop{\mathcal{E}\!\mathit{xt}}\nolimits}
\def\Sch{\mathit{Sch}}
\def\Mor{\mathop{\mathrm{Mor}}\nolimits}
\def\Ob{\mathop{\mathrm{Ob}}\nolimits}
\def\Sh{\mathop{\mathit{Sh}}\nolimits}
\def\NL{\mathop{N\!L}\nolimits}
\def\CH{\mathop{\mathrm{CH}}\nolimits}
\def\proetale{{pro\text{-}\acute{e}tale}}
\def\etale{{\acute{e}tale}}
\def\QCoh{\mathit{QCoh}}
\def\Ker{\mathop{\mathrm{Ker}}}
\def\Im{\mathop{\mathrm{Im}}}
\def\Coker{\mathop{\mathrm{Coker}}}
\def\Coim{\mathop{\mathrm{Coim}}}

% Boxtimes
%
\DeclareMathSymbol{\boxtimes}{\mathbin}{AMSa}{"02}

%
% Macros for moduli stacks/spaces
%
\def\QCohstack{\mathcal{QC}\!\mathit{oh}}
\def\Cohstack{\mathcal{C}\!\mathit{oh}}
\def\Spacesstack{\mathcal{S}\!\mathit{paces}}
\def\Quotfunctor{\mathrm{Quot}}
\def\Hilbfunctor{\mathrm{Hilb}}
\def\Curvesstack{\mathcal{C}\!\mathit{urves}}
\def\Polarizedstack{\mathcal{P}\!\mathit{olarized}}
\def\Complexesstack{\mathcal{C}\!\mathit{omplexes}}
% \Pic is the operator that assigns to X its picard group, usage \Pic(X)
% \Picardstack_{X/B} denotes the Picard stack of X over B
% \Picardfunctor_{X/B} denotes the Picard functor of X over B
\def\Pic{\mathop{\mathrm{Pic}}\nolimits}
\def\Picardstack{\mathcal{P}\!\mathit{ic}}
\def\Picardfunctor{\mathrm{Pic}}
\def\Deformationcategory{\mathcal{D}\!\mathit{ef}}
\usepackage{cleveref}
\crefname{originallemma}{original lemma}{original lemmas}
\Crefname{originallemma}{Original lemma}{Original lemmas}
\makeatletter
\newlabel{algebra-lemma-regular-normal}{{0567}{source}{}{sourcepointer-1}{}}
\newlabel{algebra-lemma-regular-normal@cref}{{[originallemma][0][]0567}{[1][0][]source}{}{}{}}
\makeatother
\crefname{originalreference}{original reference}{original references}
\Crefname{originalreference}{Original reference}{Original references}
\makeatletter
\newlabel{lemma-blowup-regular}{{0AGR}{source}{}{sourcepointer-2}{}}
\newlabel{lemma-blowup-regular@cref}{{[originalreference][0][]0AGR}{[1][0][]source}{}{}{}}
\makeatother
\crefname{originalreference}{original reference}{original references}
\Crefname{originalreference}{Original reference}{Original references}
\makeatletter
\newlabel{lemma-cohomology-of-blowup}{{0AGS}{source}{}{sourcepointer-3}{}}
\newlabel{lemma-cohomology-of-blowup@cref}{{[originalreference][0][]0AGS}{[1][0][]source}{}{}{}}
\makeatother
\crefname{originalreference}{original reference}{original references}
\Crefname{originalreference}{Original reference}{Original references}
\makeatletter
\newlabel{lemma-dominate-by-blowing-up-in-points}{{0AHI}{source}{}{sourcepointer-4}{}}
\newlabel{lemma-dominate-by-blowing-up-in-points@cref}{{[originalreference][0][]0AHI}{[1][0][]source}{}{}{}}
\makeatother
\crefname{originalreference}{original reference}{original references}
\Crefname{originalreference}{Original reference}{Original references}
\makeatletter
\newlabel{lemma-exact-sequence}{{0AXF}{source}{}{sourcepointer-5}{}}
\newlabel{lemma-exact-sequence@cref}{{[originalreference][0][]0AXF}{[1][0][]source}{}{}{}}
\makeatother
\title{Regular local surfaces define rational singularities}
\author{The Stacks Project Authors}
\date{stacks-project:resolve.tex:local-snapshot-20261012; GNU FDL 1.2 or later}
\begin{document}
\maketitle
\noindent\textit{Editorial note.} All delivered mathematical text is the author's own, copied from the pinned Stacks Project chapter file resolve.tex; no statement or proof step was rewritten, reordered or removed.\par\medskip
\noindent\textit{Editorial note.} The counted proof is the author's complete argument: the parenthetical remark that the Nagata hypothesis is only used for the definition, the domination by a chain of blowups in closed points, and the reduction to the cohomology computation for a single blowup.\par\medskip
\noindent\textit{Editorial note.} The label lemma-regular-rational names the counted statement, and the proof environment immediately following it in the source is delivered verbatim as its proof.\par\medskip
\noindent\textit{Editorial note.} The counted statement is carried by the author's own lemma environment in the chapter source.\par\medskip
\noindent\textit{Editorial note.} Pre-open gate: python operations/agent-authoring/tools/precard\_gate.py --source resolve.tex --label lemma-regular-rational --start 2204 --end 2230 --local-id 28108 returned RESULT=OK (L28108 free, occupancy FREE, public interval clear). Pre-assembly tag self-check: canonical tag resolved and not present in reference/public-items.jsonl.\par\medskip
\noindent\textit{Editorial note.} The statement and proof are copied from the pinned chapter file; the cited inputs are kept as external pointers with their canonical Stacks tags.\par\medskip
\noindent\textit{Allowed standard prerequisites (editorial).}\begin{itemize}
\item Dominating a modification by a sequence of blowups in closed points. Same chapter, the lemma producing a sequence of blowups in closed points dominating a given modification. Provides a regular modification through which the given modification factors, reducing the vanishing to the blowup case.
\item Regularity of the blowups in the sequence. Same chapter, the lemma that the blowup of the maximal ideal of a regular local surface is regular. Shows that all the blowups in the sequence are normal, so that the comparison lemma applies.
\item The exact sequence comparing first cohomology along a modification and that of the exceptional sheaf. Same chapter, the lemma giving the short exact sequence comparing the first cohomology of a modification with that of its blowup, and the lemma computing the cohomology of the blowup. Injects the first cohomology of the modification into that of the last blowup, and reduces the remaining vanishing to the vanishing of the first direct image of each single blowup.
\end{itemize}\medskip
\ifdefined\hypertarget\hypertarget{sourcepointer-1}{}\fi
\noindent\textit{Source reference.} Original lemma 0567: sources/stacks-chow/tags.txt, source lines 6700--6700. Prerequisite pointer invoked at proof line(s) 2222; it names a standard fact of a sibling chapter that the counted argument uses, and every step of the counted proof built on it is retained verbatim in the delivered body.\par\medskip
\ifdefined\hypertarget\hypertarget{sourcepointer-2}{}\fi
\noindent\textit{Source reference.} Original reference 0AGR: sources/stacks-resolve/pinned-original-source.tex, source lines 727--727. Prerequisite pointer invoked at proof line(s) 2220; it names a standard fact of the same chapter that the counted argument uses, and every step of the counted proof built on it is retained verbatim in the delivered body.\par\medskip
\ifdefined\hypertarget\hypertarget{sourcepointer-3}{}\fi
\noindent\textit{Source reference.} Original reference 0AGS: sources/stacks-resolve/pinned-original-source.tex, source lines 786--786. Prerequisite pointer invoked at proof line(s) 2229; it names a standard fact of the same chapter that the counted argument uses, and every step of the counted proof built on it is retained verbatim in the delivered body.\par\medskip
\ifdefined\hypertarget\hypertarget{sourcepointer-4}{}\fi
\noindent\textit{Source reference.} Original reference 0AHI: sources/stacks-resolve/pinned-original-source.tex, source lines 1146--1146. Prerequisite pointer invoked at proof line(s) 2215; it names a standard fact of the same chapter that the counted argument uses, and every step of the counted proof built on it is retained verbatim in the delivered body.\par\medskip
\ifdefined\hypertarget\hypertarget{sourcepointer-5}{}\fi
\noindent\textit{Source reference.} Original reference 0AXF: sources/stacks-resolve/pinned-original-source.tex, source lines 1901--1901. Prerequisite pointer invoked at proof line(s) 2223; it names a standard fact of the same chapter that the counted argument uses, and every step of the counted proof built on it is retained verbatim in the delivered body.\par\medskip
\begin{lemma}
\label{lemma-regular-rational}
Let $A$ be a Nagata regular local ring of dimension $2$. Then $A$ defines
a rational singularity.
\end{lemma}

\begin{proof}
(The assumption that $A$ be Nagata is not necessary for this proof,
but we've only defined the notion of rational singularity in the
case of Nagata $2$-dimensional normal local domains.)
Let $X \to \Spec(A)$ be a modification with $X$ normal. By
Lemma \ref{lemma-dominate-by-blowing-up-in-points}
we can dominate $X$ by a scheme $X_n$ which is the last in a sequence
$$
X_n \to X_{n - 1} \to \ldots \to X_1 \to X_0 = \Spec(A)
$$
of blowing ups in closed points. By Lemma \ref{lemma-blowup-regular}
the schemes $X_i$ are regular, in particular
normal (Algebra, Lemma \ref{algebra-lemma-regular-normal}).
By Lemma \ref{lemma-exact-sequence} we have
$H^1(X, \mathcal{O}_X) \subset H^1(X_n, \mathcal{O}_{X_n})$.
Thus it suffices to prove $H^1(X_n, \mathcal{O}_{X_n}) = 0$.
Using Lemma \ref{lemma-exact-sequence} again, we
see that it suffices to prove $R^1(X_i \to X_{i - 1})_*\mathcal{O}_{X_i} = 0$
for $i = 1, \ldots, n$. This follows from
Lemma \ref{lemma-cohomology-of-blowup}.
\end{proof}

\end{document}
